\chapter{Night Sweats}

\medskip
\noindent\textit{\textbf{Under review.} This chapter is an AI-assisted educational draft; it is not a substitute for professional medical advice.}

\section*{Definition and Overview}
episodes of excessive sweating during sleep that soak clothing or bedding and are not explained by a warm environment. Individual assessment is important because symptoms and treatment vary with the cause.

\section*{Clinical Features}
Recurrent drenching sweats, fever, weight loss, cough, fatigue, hormonal symptoms, or anxiety, depending on the cause.

\section*{When to Seek Medical Care}
Seek medical review for new, persistent, worsening, or function-limiting symptoms. Seek emergency help for severe breathing difficulty, collapse, sudden neurological change, severe chest or abdominal pain, or any rapidly worsening illness.

\section*{Diagnosis and Investigations}
History, examination, medication and alcohol review, and targeted blood tests or imaging guided by associated symptoms rather than indiscriminate testing.

\section*{Management}
Treat the underlying cause, improve bedroom temperature and sleep comfort, and seek review for persistent drenching sweats or associated fever, weight loss, or swollen glands.

\section*{Complications}
Possible complications depend on the cause and may include progression of the condition, effects on other organs, treatment adverse effects, and reduced quality of life. Early assessment can reduce preventable harm.

\section*{Prognosis and Outcomes}
Outcome depends on the underlying cause, severity, complications, and response to treatment. Discuss individual prognosis with the treating clinician.

\section*{Prevention and Self-Care}
Follow the treatment plan, keep relevant appointments, avoid known triggers, protect affected areas from injury, and seek help early if symptoms worsen. Do not start or stop prescription treatment without clinical advice.

\section*{Expanded clinical context}
Night sweats is a clinical condition. Interpretation should account for age, baseline health, medicines, comorbidities, goals, and access to care. This chapter supports discussion with a qualified clinician and does not establish a diagnosis.

\section*{Assessment and decision points}
Assessment may combine history and examination with targeted blood tests, imaging, monitoring, or specialist review. Test results require clinical context. Document the main concern, time course, functional impact, relevant risks, and red flags. Use targeted investigations when their result is likely to change management.

\section*{Management and follow-up}
Management may include education, observation, lifestyle measures, medicines, rehabilitation, procedures, or referral, with benefits and risks discussed. Explain expected improvement, when reassessment is needed, and how follow-up can be accessed. Avoid starting, stopping, or sharing prescription medicines without professional advice.

\section*{Safety-netting}
Seek urgent care for sudden severe symptoms, chest pain, breathing difficulty, collapse, new neurological deficit, or rapidly worsening function.

\section*{Practical prevention}
Follow the care plan, keep a current medicine list, attend monitoring, and arrange review for persistent or changing symptoms. Keep written instructions and seek review if the topic recurs, changes, or does not improve as expected.

\section*{Limitations}
This educational expansion should be checked against current local guidance and authoritative topic-specific sources.
\clearpage
